
\documentclass[a4paper, 12]{article}


\usepackage{amsmath,amssymb} 
\usepackage[french]{babel}
                                                                                                                               
\newtheorem{ex}{Exemple}

\newcommand\set[1]{\ensuremath{\{#1\}}}


\begin{document}

\title{Qualifier les solutions d’un énoncé de \emph{Kimind}}
\date{}
\maketitle

\section{Introduction}
Ce document explore les caractéristiques des solutions associées aux énoncés du jeu \emph{Kimind}, on distingue les solutions \emph{dichotomiques}, \emph{aux valeurs limites}, \emph{dans l’ordre des couleurs}, \emph{l’ordre inverse des couleurs}.

\section{Solution d’un énoncé \emph{Kimind}}
Un énoncé est complètement défini par une somme et un ordre des couleurs, par exemple  \emph{instance(22, r$<$b$<$g$<$y)}. Une solution d'un énoncé est un ensemble de cartes satisfaisant les contraintes de l'énoncé sans tenir compte des positions, par exemple  \emph{$\{(red,2),$ $(blue,5),\ $ $ (green,6),$ $(yellow,9)\}$}. 

\section{Solution dichotomique}
Une solution est \emph{dichotomique} lorsque ses valeurs sont proches du milieux des domaines de valeurs imposés par les contraintes de l’énoncé. Considérons l’énoncé \emph{instance(22, r$<$b$<$g$<$y)}, il impose les domaines suivants pour les cartes de chaque couleurs : 
\set{red(1, 4, 2.5), blue(2, 6, 4), green(5, 9, 7), yellow(7, 10, 8.5)} où $red(1, 4, 2.5)$ signale que le domaine de la carte de couleur \emph{red} est \set{1,\cdots,4} et que le milieu du domaine est $2.5$.

Considérons la solution   \emph{$\{(red,2),$ $(blue,5),\ $ $ (green,6),$ $(yellow,9)\}$} associée à l’énoncé et sommons les distances relatives de chaque valeur au milieu de son domaine : 
$$ d(s) = \sum_{i=1}^{n} \frac{|v_i - m_i|}{d_i},$$ 
où $v_i$ est la valeur de la $i^{eme}$ carte de la solution, $m_i = \frac{min_i + max_i}{2}$ le milieu , $max_i$ la plus grande valeur et $min_i$ la plus petite valeur, du domaine de cette carte,  et $d_i = \frac{max_i - min_i}{2}$ est la distance du $m_i$ à la borne inférieure $min_i$. Par convention $d(s) = 1$ (plus grande distance relative possible) si 
$max_i = min_i$ afin de défavoriser une solution dont des domaines seraient des singletons.
On observe que $d(s)$ est d’autant plus petit que les valeurs sont proches des milieux des domaines et que les domaines sont grands.

Nous considérons qu’une solution $s$ est dichotomique lorsqu’elle fait partie des solutions de l’énoncé minimisant cette distance. Pour cela nous énumérons toutes les solutions de l’énoncé et calculons $d(s)$ pour chacune d’elle. Nous obtenons :

\begin{verbatim}
  d({c(blue,4),c(green,7),c(red,2),c(yellow,9)} = 1 ,
  d({c(blue,4),c(green,7),c(red,3),c(yellow,8)}) =1,
  d({c(blue,4),c(green,6),c(red,3),c(yellow,9)}) = 2,
  d([c(blue,5),c(green,7),c(red,2),c(yellow,8)}) = 2,
  d({c(blue,3),c(green,7),c(red,2),c(yellow,10)}) = 3,
  d({c(blue,3),c(green,8),c(red,2),c(yellow,9)}) = 3,
  d({c(blue,4),c(green,6),c(red,2),c(yellow,10)}) = 3,
  d({c(blue,4),c(green,7),c(red,1),c(yellow,10)}) =3,
  d({c(blue,4),c(green,8),c(red,1),c(yellow,9)}) = 3,
  d({c(blue,5),c(green,6),c(red,2),c(yellow,9)}) = 3,
  d({c(blue,5),c(green,6),c(red,3),c(yellow,8)}) = 3,
  d({c(blue,5),c(green,7),c(red,1),c(yellow,9)}) = 3,
  d({c(blue,4),c(green,5),c(red,3),c(yellow,10)}) = 4,
  d({c(blue,6),c(green,7),c(red,1),c(yellow,8)}) = 4,
  d({c(blue,3),c(green,8),c(red,1),c(yellow,10)}) = 5 ,
  d({c(blue,5),c(green,6),c(red,1),c(yellow,10)}) = 5,
  d({c(blue,5),c(green,6),c(red,4),c(yellow,7)}) = 5,
  d({c(blue,2),c(green,9),c(red,1),c(yellow,10)}) = 7
\end{verbatim}

On forme une série statistique $S$ où chaque élément est de la forme $d(s, d(s))$ où $sol$ est la solution et $d(s)$ la distance.

Nous constituons une nouvelle série statistique $D$ où chaque élément est \emph{e(d, n)} où $d$ est une distance et \emph{n}  le nombre de fois où elle apparaît dans $S$. Ainsi 
$D = \set{e(1, 2), e(2, 2), e(3, 8), e(4, 2), e(5, 3), e(7, 1) }$.

Nous partitionnons $D$ avec des quintiles $D = \set{D_1, D_2, D_3, D_4, D_5}$. Une solution $s$ sera qualifiée de dichotomique si sa distance $d$ se trouve dans $D_1$, donc parmi les plus petites distances. Ainsi les distances précédentes sont partitionnées commet suit :
\begin{verbatim}

  D1 = [e(1.0,2),e(2.0,2)],
  D2 = [e(3.0, 8)],
  D3 = [],
  D4 = [e(4.0,2),e(5.0,3)],
  D5 = [e(7.0,1)]
\end{verbatim}

La distance de solution  \emph{$\{(red,2),$ $(blue,5),\ $ $ (green,6),$ $(yellow,9)\}$} valant $3$, elle n’est pas dans dans $D_1$ aussi cette solution n’est pas considérée comme propice à un traitement dichotomique.

\section{Solution limite}
Une solution est \emph{limite} lorsque ses valeurs sont proches des limites des domaines de valeurs imposés par les contraintes de l’énoncé. 

On procède de la même façon qu’avec les solutions dichotomiques en calculant pour chaque valeur d’une solution $s$ la somme  de sa distance relative à la plus proche limite de son domaine :
$$ l(s) = \sum_{i=1}^{n} \frac{min(v_i - min_i, max_i - v_i)}{d_i},$$ 
où $v_i$ est la valeur de la $i^{eme}$ carte de la solution, $max_i$ la plus grande valeur et $min_i$ la plus petite valeur, du domaine de cette carte,  et $d_i = \frac{max_i - min_i}{2}$ est la distance du milieu du domaine  à la borne inférieure $min_i$ c’est à dire la plus grande distance possible par rapport à une borne du domaine. Par convention $l(s) = 0$ si 
$max_i = min_i$).
On observe que $l(s)$ est d’autant plus petit que les valeurs sont proches des limites des domaines et que les domaines sont grands.

Nous considérons qu’une solution $s$ est limite lorsqu’elle fait partie des solutions de l’énoncé minimisant cette distance. Pour cela nous énumérons toutes les solutions de l’énoncé et calculons $l(s)$ pour chacune d’elle. 

On forme une série statistique $S$ où chaque élément est de la forme $d(s, l(s))$ où $sol$ est la solution et $l(s)$ la distance.

Nous constituons une nouvelle série statistique $D$ où chaque élément est \emph{e(l, n)} où $l$ est une distance et \emph{n}  le nombre de fois où elle apparaît dans $S$. Ainsi 
$D = \set{e(1, 2), e(2, 2), e(3, 8), e(4, 2), e(5, 3), e(7, 1) }$.

Nous partitionnons $D$ avec des quintiles $D = \set{D_1, D_2, D_3, D_4, D_5}$. Une solution $s$ sera qualifiée de limite si sa distance $l(s)$ se trouve dans $D_1$, donc parmi les plus petites distances. 

\end{document}